\section{Out-of-Sample Evidence and Model Risk}

\noindent\begin{minipage}{\columnwidth}
{\centering
  \includegraphics[width=0.95\columnwidth]{\FigOosRmse}
  \captionof{figure}{Rolling repricing on \OosDates{} origins and \OosObservations{} common targets, $q=0.40\%$, same search policy for every model. The dashed line is the error of IV-surface persistence on its support.}
  \label{fig:oos_r2}
\par}
\end{minipage}\par

\noindent\begin{minipage}{\columnwidth}
{\centering
  \includegraphics[width=0.95\columnwidth]{\FigOosRsquared}
  \captionof*{figure}{\textbf{Figure~\ref{fig:oos_r2} (continued).} Equal-date $R^2_{OOS}$ against the origin mean IV and against persistence; positive values indicate improvement over the benchmark.}
\par}
\end{minipage}\par

\subsection{Conditional next-date repricing}
Each origin is recalibrated on its own CC64 sample with a global search for every model; no model starts from an earlier origin. Model comparisons use \OosObservations{} targets; persistence comparisons use the \OosPersistenceObservations{} inside the origin interpolation domain, so the two benchmarks have different supports. Equal date weighting limits the influence of densely populated dates.

The mean daily IV RMSE is \OosBSRMSE{}~bp for Black--Scholes, \OosHRMSE{} for Heston, \OosBRMSE{} for Bates and \OosWRMSE{} for Bates--Hawkes. Against the origin mean IV the equal-date $R^2_{OOS}$ is \OosBSRMean\%, \OosHRMean\%, \OosBRMean\% and \OosWRMean\%. Against persistence it is \OosBSRPersistence\%, \OosHRPersistence\%, \OosBRPersistence\% and \OosWRPersistence\%: no model beats carrying the previous surface forward, and Bates--Hawkes is better than persistence on \OosWBeatsPersistence{} of \OosDates{} dates.

\begin{findingbox}[Model-risk conclusion]
Diebold--Mariano tests on daily MSE with Newey--West variance \citep{diebold1995,newey1987} rank Bates--Hawkes above Black--Scholes and Heston ($p<0.001$), but not above Bates ($p=\OosDMHawkesVsBP$) and not away from persistence ($p=\OosWDMPersistenceP$). With 30 dependent dates the richest model is the best of the four, not a demonstrated improvement on the simplest benchmark.
\end{findingbox}

\subsection{Jump-layer interpretation}
The self-exciting extension is motivated by clustered option-implied jump activity: relative to a constant Poisson rate, the Hawkes intensity permits bursts followed by endogenous decay. The fitted branching ratio of \HawkesBranching{} and half-life of \HawkesHalfLifeDays{} days are risk-neutral parameters, not evidence of physical news clustering; without a dated news filtration or parameter-stability evidence, these bursts cannot be identified as responses to geopolitical events.

\subsection{Interpretive limitations and validation}
Reproducibility, numerical accuracy and economic validity are distinct. A good option fit cannot validate the GLD-to-gold mapping or the scenario/WACC combination. The experiment identifies sensitivity to a chosen numerical contract, not corporate fair value.

No per-share value is reported: debt, cash and other claims are not reconciled. A genuine valuation must start from one dated filing and the matching listed security, establish ownership and attributable production, and distinguish component margin from reported operating profit. Depreciation, cash taxes, working capital, sustaining capital expenditure and growth investment must then enter consistently, with no deduction counted twice. Net debt, non-operating assets, streams and other claims must be matched to a diluted share count on the same basis.

\subsection{Operating uncertainty and research priorities}
Production, grade, recovery, mine mix and costs may depend on gold prices, inflation and exchange rates; operating dependence requires explicit assumptions. A mining-company terminal block needs an explicit resource interpretation: continued production may require reserve conversion, replacement investment or new projects, while depletion and closure liabilities shorten the cash-flow stream. A positive WACC-growth spread controls a denominator but does not justify a perpetual operating horizon. Market maturities end at 653 days, well before 2030 and the perpetuity, so fixed parameters beyond that horizon are an extrapolation assumption. A finite-reserve model and a bucketed term structure remain necessary for an economic valuation.

\section*{Conclusion}

Holding operations and corporate assumptions fixed isolates the effect of gold-price model choice. Nested-candidate safeguards improve the current Bates and Hawkes fits. Historical option-surface OOS remains distinct and no original engine beats interpolated persistence. Numerical and terminal sensitivity prevent a corporate value ranking.

The copper extension admits three producing mines and adds a conditional five-year operating contribution. Historical snapshots support actual temporal testing, but the sparse rolling experiment does not beat IV persistence. Its positive fixed-origin Bates and Hawkes comparisons are exploratory. The combined output remains an attributable component-margin sensitivity, without a reconciled FCFF or equity bridge.

The persistent-cost assumption is prudential only relative to otherwise identical efficiency-improvement scenarios, not a lower bound. The contribution is an auditable interdisciplinary framework, not a certified value: commodity basis, probability measure, operating uncertainty, corporate accounting and finite-resource closure remain the priorities for economic validation.


The next required result is a validated physical commodity forecast for each admitted model, with an explicit link to option calibration and Barrick's realised prices. This is a prerequisite for the intended DCF, not an optional robustness check. Until price, operating and accounting gates are met, the current calculations remain conditional sensitivities.


\begingroup
\setlength{\bibsep}{0.6pt}
\fontsize{6.75}{7.75}\selectfont
% Keep the frozen bibliography in a regular source file.  Overleaf compiles
% the main document with job name `output`, so depending on a generated .bbl
% file with a different name is fragile when the project is imported.
% Frozen bibliography for direct pdfLaTeX compilation on Overleaf.
% The unused bibliography ledger is archived outside the Overleaf project.
\begin{thebibliography}{}

\bibitem[Akaike, 1974]{akaike1974}
Akaike, H. (1974).
\newblock A new look at the statistical model identification.
\newblock {\em IEEE Transactions on Automatic Control}, 19(6):716--723.

\bibitem[Bates, 1996]{bates1996}
Bates, D.~S. (1996).
\newblock Jumps and stochastic volatility: Exchange rate processes implicit in
  deutsche mark options.
\newblock {\em The Review of Financial Studies}, 9(1):69--107.

\bibitem[Damodaran, 2012]{damodaran2012investment}
Damodaran, A. (2012).
\newblock {\em Investment Valuation: Tools and Techniques for Determining the
  Value of Any Asset}.
\newblock Wiley, 3 edition.

\bibitem[Glasserman, 2004]{glasserman2004monte}
Glasserman, P. (2004).
\newblock {\em Monte Carlo Methods in Financial Engineering}.
\newblock Springer.

\bibitem[Hawkes, 1971]{hawkes1971}
Hawkes, A.~G. (1971).
\newblock Spectra of some self-exciting and mutually exciting point processes.
\newblock {\em Biometrika}, 58(1):83--90.

\bibitem[Heston, 1993]{heston1993}
Heston, S.~L. (1993).
\newblock A closed-form solution for options with stochastic volatility with
  applications to bond and currency options.
\newblock {\em The Review of Financial Studies}, 6(2):327--343.

\bibitem[Schwarz, 1978]{schwarz1978}
Schwarz, G. (1978).
\newblock Estimating the dimension of a model.
\newblock {\em The Annals of Statistics}, 6(2):461--464.


\bibitem[Fang and Oosterlee, 2008]{cos2008}
Fang, F. and Oosterlee, C.~W. (2008). A novel pricing method for European options based on Fourier-cosine series expansions. {\em SIAM J. Sci. Comput.}, 31:826--848.
\bibitem[Svensson, 1994]{nss1994}
Svensson, L. (1994). \href{https://www.nber.org/papers/w4871}{Estimating and interpreting forward interest rates}. NBER WP 4871.
\bibitem[Barrick, 2026]{barrick2026}
Barrick (2026). \href{https://www.barrick.com/English/investors/quarterly-reports/default.aspx}{Q1 and Q2 financial statements, MD\&A and mine statistics}.
\bibitem[Cboe, 2026]{cboe2026}
Cboe (2026). \href{https://www.cboe.com/exchange-traded-stock/etp-options-spec}{ETP options contract specifications}.
\bibitem[SPDR, 2025]{spdr2025}
SPDR (2025). \href{https://files.spdrgoldshares.com/spdr-cms/2025-10/GLD-FAQs-January-2025.pdf}{GLD frequently asked questions}.
\bibitem[U.S. Treasury, 2026]{treasury2026}
U.S. Treasury (2026). \href{https://home.treasury.gov/resource-center-data-chart-center/interest-rates}{Daily Treasury par yield curve rates}, 2 September.

\bibitem[Barrick, 2025]{barrickar2025}
Barrick (2025). \href{https://s25.q4cdn.com/322814910/files/doc_financial/annual_reports/2025/Barrick_Annual_Report_2025.pdf}{Annual Report 2025}, MD\&A and joint-arrangements notes.
\bibitem[Barrick, 2026b]{barrickmine2026}
Barrick (2026). \href{https://s25.q4cdn.com/322814910/files/doc_financial/quarterly_results/2026/q2/Barrick_Q2_2026_Mine_Stats.pdf}{Q2 2026 mine statistics}, copper table, p. 8.

\bibitem[Schwartz, 1997]{schwartz1997}
Schwartz, E.~S. (1997). \href{https://doi.org/10.1111/j.1540-6261.1997.tb02721.x}{The stochastic behavior of commodity prices: implications for valuation and hedging}. {\em Journal of Finance}, 52(3):923--973.

\bibitem[Pan, 2002]{pan2002}
Pan, J. (2002). \href{https://doi.org/10.1016/S0304-405X(01)00088-5}{The jump-risk premia implicit in options: evidence from an integrated time-series study}. {\em Journal of Financial Economics}, 63(1):3--50.

\bibitem[Barrick Gold Research Teams, 2026]{sfcgithub}
Barrick Gold Research Teams. (2026).
\newblock {\em Operating-value sensitivity: code, data and audit tables}.
\newblock Starting Finance Club PoliTo, GitHub repository.
\newblock \href{https://github.com/StartingFinanceClubPoliTo/Research/tree/main/Barrick-Gold/Barrick\%20Mining\%20PAPER}{Barrick Mining research repository}, revised 5 October 2026.

\bibitem[Sch\"onbucher, 2001]{schonbucher2001}
Sch\"onbucher, P.~J. (2001), \emph{A LIBOR Market Model with Default Risk},
Bonn Econ Discussion Papers 15/2001, Theorem~1,
\href{https://www.econstor.eu/bitstream/10419/78414/1/bgse15_2001.pdf}{author manuscript}.

\bibitem[Sokol and Hansen, 2015]{sokol-hansen2015}
Sokol, A. and Hansen, N.~R. (2015), \emph{Exponential martingales and changes of
measure for counting processes}, Stochastic Analysis and Applications 33(5),
823--843, \href{https://doi.org/10.1080/07362994.2015.1040890}{doi:10.1080/07362994.2015.1040890}.

\bibitem[Lord et~al., 2010]{lord2010}
Lord, R., Koekkoek, R., and van Dijk, D. (2010).
\newblock A comparison of biased simulation schemes for stochastic volatility
  models.
\newblock {\em Quantitative Finance}, 10(2):177--194.

\bibitem[Diebold and Mariano, 1995]{diebold1995}
Diebold, F.~X. and Mariano, R.~S. (1995).
\newblock Comparing predictive accuracy.
\newblock {\em Journal of Business \& Economic Statistics}, 13(3):253--263.

\bibitem[Newey and West, 1987]{newey1987}
Newey, W.~K. and West, K.~D. (1987).
\newblock A simple, positive semi-definite, heteroskedasticity and
  autocorrelation consistent covariance matrix.
\newblock {\em Econometrica}, 55(3):703--708.

\end{thebibliography}

\endgroup
